\documentclass[12pt]{article}
\usepackage{amsmath, amssymb, amsthm}
\usepackage{geometry}
\usepackage{booktabs}
\usepackage{hyperref}
\usepackage{xcolor}
\geometry{margin=1in}

\newtheorem{definition}{Definition}
\newtheorem{theorem}{Theorem}
\newtheorem{lemma}{Lemma}

\title{\textbf{Gravity Softening in SPH: A First-Principles Derivation}\\[0.5em]
\large From Poisson's Equation to Lookup Tables}
\author{SPH Code Mathematical Documentation}
\date{\today}

\begin{document}
\maketitle

\begin{abstract}
This document provides a complete first-principles derivation of gravity softening
kernels for SPH simulations. We start from Poisson's equation, derive the kernel
convolution approach for all spatial dimensions (1D, 2D, 3D), and prove that the
Hernquist-Katz (1989) cubic spline softening is mathematically equivalent to
convolving the density field with the SPH kernel. We also cover the Wendland C4
kernel and discuss lookup table implementation for computational efficiency.
\end{abstract}

\tableofcontents
\newpage

%==============================================================================
\part{Foundations}
%==============================================================================

\section{The Gravity Problem in N-Body Simulations}

In N-body simulations, the gravitational force between particles diverges as $r \to 0$:
\begin{equation}
\mathbf{F}_{ij} = -\frac{Gm_i m_j}{r^2} \hat{\mathbf{r}}_{ij}
\end{equation}

This singularity causes:
\begin{enumerate}
\item Numerical instability when particles approach closely
\item Artificial two-body relaxation in collisionless systems
\item Spurious energy exchange in SPH simulations
\end{enumerate}

\textbf{Solution:} Soften the gravitational interaction at small scales.

\section{First Principles: Poisson's Equation}

All gravity derivations start from Poisson's equation:
\begin{equation}
\boxed{\nabla^2 \phi = 4\pi G \rho}
\end{equation}

where $\phi$ is the gravitational potential and $\rho$ is the mass density.

The gravitational force is:
\begin{equation}
\mathbf{F} = -m \nabla \phi
\end{equation}

\subsection{Green's Functions by Dimension}

The solution to Poisson's equation for a point mass $M$ at the origin differs by dimension:

\begin{table}[h]
\centering
\begin{tabular}{ccccc}
\toprule
\textbf{Dim} & \textbf{Geometry} & \textbf{Laplacian} & \textbf{Potential $\phi$} & \textbf{Force $F$} \\
\midrule
1D & Infinite slab & $\frac{d^2\phi}{dx^2}$ & $-2\pi G M |x|$ & $-2\pi GM \cdot \text{sign}(x)$ \\[4pt]
2D & Infinite cylinder & $\frac{1}{r}\frac{d}{dr}\left(r\frac{d\phi}{dr}\right)$ & $-2GM \ln r$ & $-\frac{2GM}{r}$ \\[4pt]
3D & Point mass & $\frac{1}{r^2}\frac{d}{dr}\left(r^2\frac{d\phi}{dr}\right)$ & $-\frac{GM}{r}$ & $-\frac{GM}{r^2}$ \\
\bottomrule
\end{tabular}
\caption{Green's functions for Poisson's equation by spatial dimension}
\end{table}

\section{The SPH Kernel}

In SPH, a particle's mass is smoothed over space using a kernel function $W$:
\begin{equation}
\rho(\mathbf{r}) = \sum_j m_j W(|\mathbf{r} - \mathbf{r}_j|, h)
\end{equation}

\subsection{Cubic Spline Kernel (M4)}

The standard cubic spline kernel in 3D is:
\begin{equation}
W(r, h) = \frac{\sigma}{h^3} \times w(q), \quad q = \frac{r}{h}
\end{equation}

where $\sigma = 1/\pi$ in 3D, and:
\begin{equation}
\boxed{w(q) = \begin{cases}
1 - \frac{3}{2}q^2 + \frac{3}{4}q^3 & 0 \leq q < 1 \\[4pt]
\frac{1}{4}(2-q)^3 & 1 \leq q < 2 \\[4pt]
0 & q \geq 2
\end{cases}}
\end{equation}

\textbf{Key property:} Compact support with $q \in [0, 2]$ (i.e., $r \in [0, 2h]$).

%==============================================================================
\part{3D Gravity: Kernel Convolution Derivation}
%==============================================================================

\section{Setup: Smoothed Density from Single Particle}

Consider a single particle of mass $m$ at the origin. Its smoothed density is:
\begin{equation}
\rho(r) = m W(r, h) = \frac{m}{\pi h^3} w(r/h)
\end{equation}

We want to find the gravitational potential $\phi(r)$ by solving Poisson's equation.

\section{Step 1: Poisson's Equation in Spherical Coordinates}

For spherically symmetric $\rho(r)$:
\begin{equation}
\frac{1}{r^2} \frac{d}{dr}\left(r^2 \frac{d\phi}{dr}\right) = 4\pi G \rho(r)
\end{equation}

Let $u(r) = r^2 \frac{d\phi}{dr}$. Then:
\begin{equation}
\frac{du}{dr} = 4\pi G \rho(r) r^2
\end{equation}

Integrating from 0 to $r$:
\begin{equation}
r^2 \frac{d\phi}{dr} = 4\pi G \int_0^r \rho(r') r'^2 \, dr' = G M_{\text{enc}}(r)
\end{equation}

\begin{definition}[Enclosed Mass]
\begin{equation}
\boxed{M_{\text{enc}}(r) = 4\pi \int_0^r \rho(r') r'^2 \, dr'}
\end{equation}
\end{definition}

\section{Step 2: Compute Enclosed Mass for Cubic Spline}

Substituting $\rho(r) = \frac{m}{\pi h^3} w(r/h)$:
\begin{equation}
M_{\text{enc}}(r) = 4\pi \int_0^r \frac{m}{\pi h^3} w(r'/h) r'^2 \, dr' = \frac{4m}{h^3} \int_0^r w(r'/h) r'^2 \, dr'
\end{equation}

Change variables: $q' = r'/h$, $dr' = h \, dq'$, $r'^2 = h^2 q'^2$:
\begin{equation}
M_{\text{enc}}(q) = \frac{4m}{h^3} \cdot h^3 \int_0^q w(q') q'^2 \, dq' = 4m \int_0^q w(q') q'^2 \, dq'
\end{equation}

\subsection{Region 1: $0 \leq q \leq 1$}

\begin{align}
M_{\text{enc}}(q) &= 4m \int_0^q \left(1 - \frac{3}{2}q'^2 + \frac{3}{4}q'^3\right) q'^2 \, dq' \\
&= 4m \int_0^q \left(q'^2 - \frac{3}{2}q'^4 + \frac{3}{4}q'^5\right) dq' \\
&= 4m \left[\frac{q'^3}{3} - \frac{3q'^5}{10} + \frac{q'^6}{8}\right]_0^q \\
&= 4m \left(\frac{q^3}{3} - \frac{3q^5}{10} + \frac{q^6}{8}\right)
\end{align}

\begin{equation}
\boxed{M_{\text{enc}}(q) = m \left(\frac{4q^3}{3} - \frac{6q^5}{5} + \frac{q^6}{2}\right), \quad 0 \leq q \leq 1}
\end{equation}

\textbf{Check:} At $q = 1$:
\begin{equation}
M_{\text{enc}}(1) = m\left(\frac{4}{3} - \frac{6}{5} + \frac{1}{2}\right) = m\left(\frac{40 - 36 + 15}{30}\right) = \frac{19m}{30}
\end{equation}

\subsection{Region 2: $1 \leq q \leq 2$}

\begin{equation}
M_{\text{enc}}(q) = M_{\text{enc}}(1) + 4m \int_1^q \frac{1}{4}(2-q')^3 q'^2 \, dq'
\end{equation}

Substitute $t = 2 - q'$, $q' = 2 - t$, $dq' = -dt$:
\begin{align}
\int_1^q (2-q')^3 q'^2 \, dq' &= \int_{2-q}^1 t^3 (2-t)^2 \, dt \\
&= \int_{2-q}^1 t^3 (4 - 4t + t^2) \, dt \\
&= \int_{2-q}^1 (4t^3 - 4t^4 + t^5) \, dt \\
&= \left[t^4 - \frac{4t^5}{5} + \frac{t^6}{6}\right]_{2-q}^1
\end{align}

At $t = 1$: $1 - \frac{4}{5} + \frac{1}{6} = \frac{11}{30}$

Let $s = 2 - q$:
\begin{equation}
M_{\text{enc}}(q) = \frac{19m}{30} + m\left(\frac{11}{30} - s^4 + \frac{4s^5}{5} - \frac{s^6}{6}\right)
\end{equation}

\begin{equation}
\boxed{M_{\text{enc}}(q) = m\left(1 - s^4 + \frac{4s^5}{5} - \frac{s^6}{6}\right), \quad s = 2-q, \quad 1 \leq q \leq 2}
\end{equation}

\textbf{Check:} At $q = 2$ ($s = 0$): $M_{\text{enc}}(2) = m$ \checkmark

\subsection{Region 3: $q \geq 2$}
\begin{equation}
M_{\text{enc}}(q) = m \quad \text{(all mass enclosed)}
\end{equation}

\section{Step 3: Derive the Force Kernel $g(r, h)$}

The gravitational force on a test mass $m_i$ is:
\begin{equation}
\mathbf{F} = -m_i \nabla \phi = -\frac{G m_i M_{\text{enc}}(r)}{r^2} \hat{\mathbf{r}}
\end{equation}

\begin{definition}[Force Softening Kernel]
We define $g(r, h)$ such that:
\begin{equation}
\mathbf{F}_{ij} = -G m_i m_j \cdot g(r_{ij}, h) \cdot \mathbf{r}_{ij}
\end{equation}
\end{definition}

Comparing with Newton's form $F = GM_{\text{enc}}/r^2$:
\begin{equation}
\boxed{g(r, h) = \frac{M_{\text{enc}}(r)}{m \cdot r^3} = \frac{M_{\text{enc}}(q)}{m \cdot h^3 q^3}}
\end{equation}

\subsection{Force Kernel: Region 1 ($0 \leq q \leq 1$)}
\begin{align}
g(q) &= \frac{1}{h^3 q^3} \left(\frac{4q^3}{3} - \frac{6q^5}{5} + \frac{q^6}{2}\right) \\
&= \frac{1}{h^3} \left(\frac{4}{3} - \frac{6q^2}{5} + \frac{q^3}{2}\right)
\end{align}

\begin{equation}
\boxed{g(q) = \frac{1}{h^3}\left(\frac{4}{3} - \frac{6q^2}{5} + \frac{q^3}{2}\right), \quad 0 \leq q \leq 1}
\end{equation}

\subsection{Force Kernel: Region 2 ($1 \leq q \leq 2$)}
\begin{equation}
\boxed{g(q) = \frac{1}{h^3 q^3}\left(1 - s^4 + \frac{4s^5}{5} - \frac{s^6}{6}\right), \quad s = 2-q}
\end{equation}

\subsection{Force Kernel: Region 3 ($q \geq 2$)}
\begin{equation}
\boxed{g(q) = \frac{1}{r^3} \quad \text{(point mass)}}
\end{equation}

\section{Step 4: Derive the Potential Kernel $f(r, h)$}

The potential is obtained by integrating the force:
\begin{equation}
\phi(r) = -\int_\infty^r \frac{G M_{\text{enc}}(r')}{r'^2} \, dr'
\end{equation}

\begin{definition}[Potential Softening Kernel]
We define $f(r, h)$ such that:
\begin{equation}
\phi_{ij} = -G m_j \cdot f(r_{ij}, h)
\end{equation}
\end{definition}

For $r \geq 2h$ (point mass region):
\begin{equation}
f(r) = \int_r^\infty \frac{1}{r'^2} \, dr' = \frac{1}{r}
\end{equation}

For $r < 2h$, the integral requires splitting across regions. The result (after
lengthy calculation) yields the same polynomials as Hernquist-Katz derived.

%==============================================================================
\part{Hernquist-Katz (1989) Formulation}
%==============================================================================

\section{Historical Approach}

Hernquist \& Katz (1989) derived their softening kernel using a polynomial ansatz
with boundary conditions, rather than explicit kernel convolution. We show their
result is equivalent.

\subsection{Parameterization}
\begin{align}
\epsilon &= h/2 \quad \text{(softening length)} \\
u &= r/\epsilon = 2r/h \quad \text{(normalized distance)}
\end{align}

Support: $u \in [0, 2]$, corresponding to $r \in [0, h]$.

\subsection{Potential Kernel}
\begin{equation}
f(u) = \frac{1}{\epsilon} \times
\begin{cases}
-\frac{u^2}{2}\left(\frac{1}{3} - \frac{3u^2}{20} + \frac{u^3}{20}\right) + \frac{7}{5} & 0 \leq u < 1 \\[8pt]
-\frac{1}{15u} - u^2\left(\frac{4}{3} - u + \frac{3u^2}{10} - \frac{u^3}{30}\right) + \frac{8}{5} & 1 \leq u < 2 \\[8pt]
\frac{1}{u} & u \geq 2
\end{cases}
\end{equation}

\subsection{Force Kernel}
\begin{equation}
g(u) = \frac{1}{\epsilon^3} \times
\begin{cases}
\frac{4}{3} - \frac{6u^2}{5} + \frac{u^3}{2} & 0 \leq u < 1 \\[8pt]
\frac{1}{u^3}\left(-\frac{1}{15} + \frac{8u^3}{3} - 3u^4 + \frac{6u^5}{5} - \frac{u^6}{6}\right) & 1 \leq u < 2 \\[8pt]
\frac{1}{u^3} & u \geq 2
\end{cases}
\end{equation}

\subsection{Boundary Conditions}
\begin{align}
f(0) &= \frac{7}{5\epsilon} = \frac{14}{5h} \approx \frac{2.8}{h} \\
g(0) &= \frac{4}{3\epsilon^3} = \frac{32}{3h^3} \\
\lim_{u \to \infty} f(u) &= \frac{1}{r} \quad \text{(point mass)} \\
\lim_{u \to \infty} g(u) &= \frac{1}{r^3} \quad \text{(point mass)}
\end{align}

%==============================================================================
\part{Equivalence Proof}
%==============================================================================

\section{Comparing the Two Approaches}

\begin{theorem}[Equivalence]
The Hernquist-Katz force kernel with $\epsilon = h$ is mathematically equivalent
to kernel-convolved gravity using the cubic spline SPH kernel.
\end{theorem}

\begin{proof}
The kernel convolution gives (for $0 \leq q \leq 1$):
\begin{equation}
g_{\text{conv}}(q) = \frac{1}{h^3}\left(\frac{4}{3} - \frac{6q^2}{5} + \frac{q^3}{2}\right)
\end{equation}

The H-K kernel with $\epsilon = h$ (so $u = r/h = q$) gives:
\begin{equation}
g_{\text{HK}}(u) = \frac{1}{h^3}\left(\frac{4}{3} - \frac{6u^2}{5} + \frac{u^3}{2}\right)
\end{equation}

These are identical. \qed
\end{proof}

\section{The Parameter Convention Difference}

The standard H-K convention uses $\epsilon = h/2$, giving $u = 2q$. This compresses
the kernel into half the radius:

\begin{table}[h]
\centering
\begin{tabular}{lcc}
\toprule
\textbf{Property} & \textbf{Kernel Convolution} & \textbf{H-K ($\epsilon = h/2$)} \\
\midrule
Support radius & $2h$ & $h$ \\
Central force $g(0)$ & $\frac{4}{3h^3}$ & $\frac{32}{3h^3}$ \\
Relative strength & 1× & 8× \\
\bottomrule
\end{tabular}
\end{table}

\textbf{Key insight:} H-K is 8× stronger at center because $\epsilon^3 = h^3/8$.

%==============================================================================
\part{Lower Dimensions: 1D and 2D}
%==============================================================================

\section{1D Gravity: Infinite Slab}

\subsection{Physics}
For a slab with surface density $\Sigma$ (mass per area), Gauss's law gives:
\begin{equation}
g(x) = -2\pi G \cdot \text{sign}(x) \cdot \Sigma_{\text{enc}}(|x|)
\end{equation}

For kernel-convolved SPH gravity, we smooth the mass along $x$.

\subsection{Cumulative Distribution Function}
\begin{definition}
\begin{equation}
F(q) = \int_{-\infty}^{q} W_{1D}(q') \, dq'
\end{equation}
\end{definition}

This gives the fraction of mass ``to the left'' of position $q$.

\subsection{1D Cubic Spline Kernel}
\begin{equation}
W_{1D}(q) = \frac{2}{3} \times
\begin{cases}
1 - \frac{3}{2}q^2 + \frac{3}{4}|q|^3 & |q| \leq 1 \\[4pt]
\frac{1}{4}(2-|q|)^3 & 1 < |q| \leq 2 \\[4pt]
0 & |q| > 2
\end{cases}
\end{equation}

\subsection{Cumulative F(q)}
After integration:
\begin{equation}
F(q) =
\begin{cases}
0 & q \leq -2 \\[4pt]
\frac{1}{24}(2+q)^4 & -2 < q \leq -1 \\[4pt]
\frac{1}{24} + \frac{2}{3}\left(0.6875 + q - \frac{q^3}{2} + \frac{3q^4}{16}\right) & -1 < q \leq 1 \\[4pt]
F(1) + \frac{1}{24}\left[1 - (2-q)^4\right] & 1 < q \leq 2 \\[4pt]
1 & q > 2
\end{cases}
\end{equation}

where $F(1) = \frac{1}{24} + \frac{2}{3}(1.375) \approx 0.9583$.

\subsection{1D Gravity Calculation}
\begin{equation}
\boxed{g_i = -\pi G \sum_j m_j \left[2F\left(\frac{x_i - x_j}{h_{ij}}\right) - 1\right]}
\end{equation}

The factor $2F(q) - 1 \in [-1, +1]$ represents net mass to left minus right.

\section{2D Gravity: Infinite Cylinder}

\subsection{Physics}
For a cylinder (2D disk), the potential is logarithmic:
\begin{equation}
\phi = -2GM \ln r, \quad F = -\frac{2GM}{r}
\end{equation}

\subsection{2D Cumulative Distribution}
\begin{equation}
F(q) = \int_0^q W_{2D}(q') \cdot 2\pi q' \, dq' \Big/ \int_0^\infty W_{2D}(q') \cdot 2\pi q' \, dq'
\end{equation}

After normalization:
\begin{equation}
F(q) = \frac{20}{7} \int_0^q w(q') q' \, dq'
\end{equation}

\subsection{Result}
\begin{equation}
F(q) =
\begin{cases}
0 & q \leq 0 \\[4pt]
\frac{20}{7}\left(\frac{q^2}{2} - \frac{3q^4}{8} + \frac{3q^5}{20}\right) & 0 < q \leq 1 \\[4pt]
F(1) + \frac{5}{7}\left(0.3 - \frac{(2-q)^4}{2} + \frac{(2-q)^5}{5}\right) & 1 < q \leq 2 \\[4pt]
1 & q > 2
\end{cases}
\end{equation}

where $F(1) \approx 0.786$.

\subsection{2D Gravity Calculation}
\begin{equation}
\boxed{g(r, h) = \frac{F(r/h)}{r}, \quad \mathbf{F}_i = -2G \sum_j m_j \cdot g(r_{ij}, h_{ij}) \cdot \hat{\mathbf{r}}_{ij}}
\end{equation}

\section{Why No Explicit $F_{3D}(q)$?}

In 3D, the enclosed mass formula already gives the force directly:
\begin{equation}
F = -\frac{GM_{\text{enc}}(r)}{r^2} = -Gm \cdot g(r,h) \cdot r
\end{equation}

The H-K kernels $f(r,h)$ and $g(r,h)$ \textbf{are} the result of this integration.
No separate cumulative function is needed because the spherical geometry is
simpler than 1D/2D.

%==============================================================================
\part{Wendland C4 Kernel}
%==============================================================================

\section{Wendland C4 Definition}

The Wendland C4 kernel has compact support $q \in [0, 1]$ and higher smoothness:
\begin{equation}
W_{C4}(q) = \frac{495}{32\pi h^3} (1-q)^6 (1 + 6q + \frac{35q^2}{3}), \quad q \leq 1
\end{equation}

\section{Gravitational Potential Kernel}

The potential kernel is obtained by integrating the Green's function:
\begin{equation}
\tilde{\phi}(q) = \frac{1}{h} \sum_{n=0}^{9} a_n q^n, \quad q < 1
\end{equation}

Coefficients:
\begin{align}
a_0 &= 3.4374743761 & a_5 &= -26.3399094060 \\
a_1 &= -0.0031873250 & a_6 &= -44.1079372114 \\
a_2 &= -10.2154807743 & a_7 &= 82.6543766683 \\
a_3 &= -1.1577720555 & a_8 &= -50.5921624056 \\
a_4 &= 36.1013669755 & a_9 &= 11.2232565249
\end{align}

\section{Force Kernel}

\begin{equation}
g(q) = -\frac{1}{h^3 q} \frac{d\tilde{\phi}}{dq} = \frac{1}{h^3} \times \left(-\frac{1}{q}\sum_{n=1}^{9} n a_n q^{n-1}\right)
\end{equation}

%==============================================================================
\part{Implementation}
%==============================================================================

\section{Lookup Table Design}

For computational efficiency, kernels are precomputed in lookup tables:

\begin{table}[h]
\centering
\begin{tabular}{lccc}
\toprule
\textbf{Kernel} & \textbf{Table Size} & \textbf{Range} & \textbf{Accuracy} \\
\midrule
Hernquist-Katz $f(u)$, $g(u)$ & 4096 & $u \in [0, 2]$ & $< 10^{-6}$ \\
Wendland C4 $\phi(q)$, $g(q)$ & 2048 & $q \in [0, 1]$ & $< 10^{-6}$ \\
\bottomrule
\end{tabular}
\end{table}

\section{Linear Interpolation}

For lookup at position $u$:
\begin{align}
i &= \lfloor u / \Delta u \rfloor \\
t &= (u - i \cdot \Delta u) / \Delta u \\
f(u) &\approx f_i + t \cdot (f_{i+1} - f_i)
\end{align}

Precomputed slopes: $s_i = (f_{i+1} - f_i) / \Delta u$ for $O(1)$ evaluation.

\section{Code Functions}

\begin{table}[h]
\centering
\begin{tabular}{lll}
\toprule
\textbf{Function} & \textbf{Dimension} & \textbf{Description} \\
\midrule
\texttt{f(r, h)} & 3D & H-K potential kernel \\
\texttt{g(r, h)} & 3D & H-K force kernel \\
\texttt{wendland\_c4\_phi(r, h)} & 3D & W-C4 potential kernel \\
\texttt{wendland\_c4\_g(r, h)} & 3D & W-C4 force kernel \\
\texttt{cubic\_spline\_F\_1d(q)} & 1D & Cumulative distribution \\
\texttt{kernel\_gravity\_1d(x, h)} & 1D & 1D gravity: $2F(q)-1$ \\
\texttt{cubic\_spline\_F\_2d(q)} & 2D & Enclosed mass fraction \\
\texttt{kernel\_gravity\_2d(r, h)} & 2D & 2D gravity: $F(q)/r$ \\
\bottomrule
\end{tabular}
\end{table}

%==============================================================================
\part{Summary}
%==============================================================================

\section{Main Results}

\begin{enumerate}
\item \textbf{Kernel convolution = H-K:} The Hernquist-Katz cubic spline softening
      is mathematically equivalent to convolving the density field with the
      SPH kernel and solving Poisson's equation.

\item \textbf{Parameterization:} H-K uses $\epsilon = h/2$, giving 8× stronger
      central force than kernel convolution with the same $h$.

\item \textbf{Dimension-dependent:} The Green's function differs by dimension,
      requiring different formulations for 1D (slab), 2D (cylinder), and 3D (point mass).

\item \textbf{Computational efficiency:} Lookup tables with linear interpolation
      provide $O(1)$ evaluation with $< 10^{-6}$ relative accuracy.
\end{enumerate}

%==============================================================================
\section*{References}
%==============================================================================

\begin{enumerate}
\item Hernquist, L. \& Katz, N. (1989). TREESPH: A unification of SPH with the
      hierarchical tree method. \textit{ApJS}, 70, 419.
\item Price, D.J. \& Monaghan, J.J. (2007). An energy-conserving formalism for
      adaptive gravitational force softening. \textit{MNRAS}, 374, 1347.
\item Dehnen, W. (2001). Towards optimal softening in 3D N-body codes.
      \textit{MNRAS}, 324, 273.
\item Wendland, H. (1995). Piecewise polynomial, positive definite and compactly
      supported radial functions of minimal degree. \textit{Adv. Comput. Math.}, 4, 389.
\item Monaghan, J.J. (1992). Smoothed particle hydrodynamics. \textit{ARA\&A}, 30, 543.
\end{enumerate}

\end{document}
